\documentclass[letterpaper]{article}
\usepackage[submission]{aaai2027}
\usepackage[hyphens]{url}
\usepackage{graphicx}
\usepackage{natbib}
\urlstyle{rm}
\def\UrlFont{\rm}
\usepackage{caption}
\frenchspacing
\usepackage{amsmath}
\usepackage{amssymb}
\usepackage{amsthm}
\usepackage{booktabs}
\newtheorem{definition}{Definition}
\newtheorem{lemma}{Lemma}
\newtheorem{proposition}{Proposition}
\newtheorem{observation}{Observation}
\newtheorem{theorem}{Theorem}
\newtheorem{corollary}{Corollary}
\renewcommand{\qedsymbol}{}
\usepackage{multirow}
\usepackage{cuted}
\emergencystretch=1.2em
\pdfinfo{
/TemplateVersion (2027.1)
}
\setcounter{secnumdepth}{0}

\title{WEAVE: Affordable Real Style Transfer via Wavelet-Endpoint Aligned Vector-Field Editing }
\author{Anonymous Submission}
\affiliations{}
\begin{document}
\maketitle

\begin{abstract}
Art-to-art style transfer is uniquely susceptible to a no-op shortcut: since the source is already an artwork, an identity mapping achieves near-perfect standard metrics, making style suppression competitively viable.
We calibrate evaluation with identical-image transfer (IDT), where the unchanged source defines a no-op floor any valid method must clear.
We trace this failure to Euclidean latent geometry, where layout and texture share one coordinate basis, letting the flow objective suppress style-conditioned transport.
We provide a mechanistic analysis of this style-suppression regime and show that changing coordinates breaks the deadlock.
We propose \textnormal{WEAVE} (Wavelet-Endpoint Aligned Vector-field Editing), a wavelet-decoupled rectified flow that routes style through high-frequency subbands with endpoint-only whitening-and-coloring.
On Distinct5-WikiArt, WEAVE clears the IDT floor (CLIP-S\,$=$\,0.7213, LPIPS\,$=$\,0.2868), with 903K parameters training in \textbf{3.08 minutes} on a single \textbf{RTX 3060}.
\end{abstract}

\begin{strip}
\noindent\includegraphics[width=\textwidth]{fig_distinct5_page1_summary.pdf}
\captionof{figure}{IDT-calibrated comparison on Distinct5-WikiArt.
Left: CLIP-S versus $1-\mathrm{LPIPS}$, with bubble area encoding training time.
The red points mark two WEAVE operating points.
Square markers denote training-free diffusion editors (StyleAligned, IP-Adapter).
The dashed line marks the IDT floor (0.6933); methods below it have not moved in the requested direction.
Right: target-pooled ArtFID on the same benchmark, with numbers inside the bars indicating training or response time.}
\label{fig:page1}
\end{strip}

\section{Introduction}

Art-to-art style transfer has an unusually forgiving failure mode: the source is already an artwork, so near-identity mapping can score competitively under standard metrics such as LPIPS and ArtFID.
To expose this, we calibrate evaluation with identical-image transfer (IDT): the unchanged source defines a no-op floor, and a method below that floor has not moved in the requested direction.
On Distinct5-WikiArt, under IDT calibration, SaMam achieves strong content fidelity (LPIPS 0.2434) but its CLIP-S of 0.5816 falls below the IDT floor of 0.6933---a signature of source reconstruction rather than target-direction transfer.
IDT thus distinguishes real target-direction transfer from no-op behavior: a method below the identity CLIP-S floor has not moved in the requested direction (formalized in Section~\ref{sec:idt}).

We trace this failure to the geometry of Euclidean latent flow.
In standard latent coordinates, the rectified-flow target $u_t=z_1-z_0$ mixes low-frequency layout motion with high-frequency texture motion in one shared basis.
The flow loss then rewards structure preservation more strongly than style-conditioned motion, pulling optimization toward a style-suppression regime.

This style-suppression mechanism also explains why reference-conditioned diffusion editors often trade between content collapse and style dilution: in Euclidean latent space, semantic layout and stylistic texture compete in one coordinate system.
We characterize this regime through a mechanistic frequency-dominance analysis (Section~\ref{sec:attractor}) and show that removing the coordinate change weakens the frontier, route mismatch lowers quality, and repeated style injection increases damage.

WEAVE mitigates this suppression regime by changing coordinates rather than enlarging the generator.
It applies a single-level Haar transform, weakens LL supervision, routes style queries through high-frequency bands, and applies whitening-and-coloring only at the endpoint.
In the ideal routed limit, the active supervised transport reduces to two approximately decoupled high-frequency bands, removing the dominant low-frequency conflict from both the gradient and the style-query path.

Crucially, WEAVE strips away rather than adds complexity: a subtractive audit finds routing heuristics and multi-step injection protocols empirically redundant in this geometry.
The resulting 903K-parameter backbone trains in 3.08 minutes on a single RTX 3060, generates 750 pairs in 83.7 seconds, and adapts to new styles by updating only a compact style-memory module.

\paragraph{Contributions.}
(1)~We introduce IDT-calibrated evaluation for art-to-art transfer, distinguishing real target-direction transfer from no-op behavior.
(2)~We characterize a local style-suppression regime in Euclidean latent flow through a frequency-dominance mechanistic analysis, yielding a rigid prediction-to-evidence mapping (validated in Section 4.3).
(3)~We propose \textbf{WEAVE} (Wavelet-Endpoint Aligned Vector-field Editing), a wavelet-decoupled rectified flow with high-frequency query routing and endpoint subband alignment whose active supervised transport reduces to two decoupled high-frequency bands in the ideal routed limit (Section~\ref{sec:baselock}).
(4)~We show that this geometry, together with a subtractive audit of unnecessary complexity, enables a practical real-transfer regime at a fraction of the compute cost of competing methods, making high-quality style transfer accessible on a single consumer GPU.

\section{Related Work}
\label{sec:related}

\paragraph{Domain-conditional transfer and evaluation.}
Classical exemplar transfer uses a target image at inference time~\cite{gatys2015neural,gatys2016image,huang2017adain,li2017universal}; our setting uses only a source image and a target style identity, so the task is domain-conditional translation.
The closest learned baselines are SaMST and SaMam~\cite{liu2024samst,liu2024samam}, with larger diffusion-based editors (StyleID, SDEdit~\cite{meng2021sdedit,chung2024styleid}) and training-free approaches (StyleAligned~\cite{gal2022stylealigned}) defining the high-compute end.
We pair standard metrics (CLIP-S, LPIPS~\cite{zhang2018lpips}, ArtFID~\cite{wright2022artfid}) with IDT calibration, where the unchanged source defines a no-op floor that methods must clear.

\paragraph{Wavelet geometry.}
Prior wavelet methods mainly use multi-scale decompositions as feature-processing modules~\cite{daubechies1992ten,phung2023wavelet}.
We instead use a wavelet basis as a coordinate change that separates low-frequency structure from high-frequency style motion, letting the optimizer resolve layout and texture in approximately decoupled channels.

\section{Theory and Method}
\label{sec:method}

This section develops the mechanistic insight behind WEAVE. We first calibrate evaluation with identical-image transfer (IDT) to expose the art-to-art identity shortcut (Section~\ref{sec:idt}), then diagnose why Euclidean latent flow suppresses style through a frequency-dominance analysis (Section~\ref{sec:attractor}), and finally build a wavelet-decoupled design that escapes this regime while remaining compact enough to train in minutes on a single consumer GPU.

\begin{figure*}[t!]
\centering
\includegraphics[width=\textwidth]{framework_sfm_main.png}
\caption{\textbf{WEAVE architecture} (top to bottom). \textbf{(a)} \emph{Timeline contrast}: prior arts inject style at every ODE step (causing exponential content decay), while WEAVE applies alignment \emph{only} at $t=1$. \textbf{(b)} \emph{Wavelet decomposition}: a single-level Haar transform splits $z_t$ into LL (content, blue) and LH/HL/HH (style, orange) subbands. Style queries are routed exclusively through high-frequency bands; LL receives only weak supervision ($\lambda_{LL}\!=\!0.3$) for global color stability and is blocked from style information (\textbf{$\times$}). \textbf{(c)} \emph{Network pipeline}: the velocity network predicts $\mathbf{v}_{LL}$, $\mathbf{v}_{LH}$, $\mathbf{v}_{HL}$; HH is left as an endpoint-only band. After 8-step Heun integration, an \textbf{Endpoint Alignment} module applies whitening-and-coloring (WCT) to HF subbands together with Semantic Region SWD (K-means partition + quantile matching), producing the final stylized output $\hat{z}_1$. Decoupled frequency channels prevent content (blue) and style (orange) from competing in the same gradient.}
\label{fig:framework}
\end{figure*}



\subsection{Setup}

Let $x_c \in \mathbb{R}^{H \times W \times 3}$ be a content image and $s \in \{1,\ldots,K\}$ a target style domain.
We work in the latent space of the Stable Diffusion v1.5 EMA VAE.
We write $z_0 = \mathcal{E}(x_c) \in \mathbb{R}^{4 \times 32 \times 32}$.
During training, $z_1 \sim p_s$ is a latent sampled from the target domain. Rectified flow uses the straight segment $z_t = (1-t)z_0 + tz_1,\; u_t = z_1 - z_0,\; t \sim \mathcal U([0,1])$.
At inference we use an 8-step Heun solver (Heun gave the best cost-quality tradeoff in our controls; Euler loses style, RK4 adds cost without improving the final result).
We then decode $\hat x_s = \mathcal D(\hat z_1)$.

\subsection{The Identity Transfer (IDT) Calibration}
\label{sec:idt}

Art-to-art style transfer is uniquely susceptible to a no-op shortcut: since the source is already an artwork, standard metrics reward doing nothing.
We formalize this through identical-image transfer (IDT), where the unchanged source defines a no-op floor that any valid method must clear.

\begin{definition}[Identical-Image Transfer (IDT) Floor]
\label{def:idt}
Because the source image is already stylized, $\mathrm{CLIP}(x_{\mathrm{src}},s_{\mathrm{tgt}})$ is a natural no-op baseline. A valid transfer must satisfy $\mathrm{CLIP}(x_{\mathrm{out}},s_{\mathrm{tgt}})\ge \mathrm{CLIP}(x_{\mathrm{src}},s_{\mathrm{tgt}})$; below-floor scores indicate negative transfer---the method has not moved in the requested direction.
\end{definition}

The IDT floor turns standard metrics into direction-aware measures: once the identity row is included, low LPIPS alone is not enough---a method must also clear the floor, or it is merely reconstructing the source.
On Distinct5-WikiArt, SaMam achieves strong content fidelity (LPIPS 0.2434) but its CLIP-S of 0.5816 falls below the IDT floor of 0.6933---a signature of source reconstruction rather than target-direction transfer.
IDT thus exposes an evaluation blind spot that is particularly acute in art-to-art transfer, where the source already contains rich stylistic features.

\subsection{The Frequency-Dominance Dilemma}
\label{sec:attractor}

While analyzing global convergence in highly non-convex diffusion and flow matching models remains an open challenge, we can characterize this failure mode through a mechanistic lens based on the empirical spectral properties of natural images.
If style transfer is learned directly in Euclidean latent coordinates, the natural objective is
\begin{equation}
\mathcal L_{\mathrm{euc}} = w_f\,\mathcal L_{\mathrm{FM}} + w_s\,\mathcal L_{\mathrm{sty}},
\label{eq:euc}
\end{equation}
where $\mathcal L_{\mathrm{FM}} = \mathbb E_t\|\tilde v_\theta(z_t,t,s)-u_t\|_2^2$ and $\mathcal L_{\mathrm{sty}}$ is any endpoint style-matching term.
The difficulty is structural: the straight target $u_t=z_1-z_0$ mixes low-frequency layout motion with high-frequency texture motion in the same coordinates.

\begin{definition}[The Style-Suppression Regime]
\label{def:attractor}
We characterize the ``no-op'' failure mode through a parameter set $\Theta_0$ where the velocity field largely ignores style conditioning:
$\Theta_0 = \{\theta: \tilde v_\theta(z,t,s) \approx \bar v_\theta(z,t)\}$, where $\bar v_\theta(z,t) = \mathbb E_s[\tilde v_\theta(z,t,s)]$ is the marginal velocity averaged across all styles.
Empirically, models trapped in $\Theta_0$ exhibit strong content preservation but fall below the IDT style floor.
\end{definition}

\begin{observation}[Frequency-Dominance in Gradient Signals]
\label{prop:attractor}
Natural images and their VAE latents follow a $1/f$ spectral distribution.
Consequently, the flow-matching loss $\mathcal L_{\mathrm{FM}}$ enforces massive gradient penalties on low-frequency structural deviations ($L_f^{\mathrm{low}} \gg L_f^{\mathrm{high}}$), while style metrics (like CLIP) are relatively insensitive to high-frequency textures ($L_s^{\mathrm{high}} \approx 0$).
In a shared Euclidean basis, optimization prioritizes minimizing the dominant low-frequency flow loss, heavily suppressing the style-conditioned high-frequency gradients.
The model thus naturally gravitates toward the style-agnostic regime $\Theta_0$.
\end{observation}

A highly effective remedy is to decouple the frequency channels.
By transforming the latents via discrete wavelet transform (DWT), we can isolate the dominant low-frequency structure from the high-frequency texture, preventing them from competing within the same optimization gradient---which is precisely what wavelet decoupling does.

\subsection{Wavelet-Decoupled Transport}
\label{sec:fiberbundle}

WEAVE first changes coordinates.
A single-level Haar wavelet transform decomposes the $4{\times}32{\times}32$ latent into four orthogonal subbands: a low-frequency LL band and three high-frequency bands (LH, HL, HH), each of size $4{\times}16{\times}16$.
The transform is orthogonal, so the Frobenius norm splits additively:
\begin{equation}
\|z\|_F^2 = \|\ell\|_F^2 + \|h_1\|_F^2 + \|h_2\|_F^2 + \|h_3\|_F^2.
\end{equation}
We write $\pi(z)=\ell$ for the low-frequency projection.
Due to the orthogonality of the Haar transform, modifying only the high-frequency subbands leaves the structural base coordinate $\ell$ strictly invariant.

\textbf{Empirical frequency spectrum of VAE latents.}
On Distinct5-WikiArt, the LL subband consistently carries $>\!90\%$ of the total Frobenius norm (confirming the $1/f$ structure), and the style-induced change between a content and a target-style latent is concentrated in LH/HL/HH---2--4$\times$ larger than in LL.
This separation justifies weakening LL supervision and routing style queries through the high-frequency bands without losing structural fidelity.

\textbf{Why single-level Haar.}
We select the Haar wavelet for three practical reasons.
(1)~Exact orthogonality: the four subbands partition the signal energy without overlap or redundancy, so modifying one band never aliases into another---a property that longer wavelets (e.g., Daubechies-2) approximate but do not guarantee in the discrete case.
(2)~Compact spatial support: Haar operates on $2{\times}2$ patches with simple averages and differences, introducing no boundary artifacts or ringing; in our controls, Daubechies-2 improves CLIP-S marginally but raises LPIPS from 0.2868 to 0.3288 due to filter-induced oscillations across object boundaries.
(3)~Efficiency: the forward transform costs $O(HW)$ and adds zero learnable parameters, so it imposes negligible overhead at both training and inference.
We keep a single decomposition level because the latent resolution is already compact ($4{\times}32{\times}32$); further decomposition would shrink subbands below a useful spatial scale while adding no new frequency-band separation under the $1/f$ prior.

\subsection{High-Frequency Routing}
\label{sec:baselock}

Let $\mathcal W(z_t) = (\ell_t,h_{1,t},h_{2,t},h_{3,t})$ and $\mathcal W(u_t) = (u_\ell,u_1,u_2,u_3)$.
The WEAVE model receives $(\ell_t,h_{1,t},h_{2,t},h_{3,t})$ and predicts three velocity heads $(v_\ell,v_1,v_2)$.
The $h_3$ (HH) band has no learned velocity head; it stays in the state and is stylized only at the endpoint.

The implemented training objective is
\begin{equation}
\mathcal L_{\mathrm{impl}}
=
\mathbb E_t\!\left[
\lambda_{LL}\|v_\ell-u_\ell\|_2^2
+ \|v_1-u_1\|_2^2
+ \|v_2-u_2\|_2^2
\right],
\label{eq:weave-loss}
\end{equation}
with $\lambda_{LL} = 0.3$.
WEAVE de-weights LL rather than eliminating it: some styles also shift coarse color tone, so LL remains a weak anchor.

\textbf{Backbone architecture.}
The velocity network consists of four residual blocks (width~64, $\sim$903K trainable parameters) with three separate prediction heads---one $1{\times}1$ convolution per supervised subband ($\ell$, $h_1$, $h_2$).
The $h_3$ (HH) band receives no velocity head and evolves passively along the rectified-flow trajectory.
Style conditioning is injected through a memory module: each target style $s$ owns a learned embedding $m_s \in \mathbb R^{d}$ ($d{=}64$), which is projected to keys and values for cross-attention.
The attention query is assembled from the high-frequency subbands via a lightweight linear projection $Q[h_1; h_2; h_3]$, and the resulting features are fused into the residual-stream activations.
Critically, the LL subband enters the backbone as a separate channel and interacts with style only through the shared residual blocks---never through the attention path---preserving the decoupled gradient structure described in Section~\ref{sec:attractor}.

Loss-level de-weighting is only half of the design.
At inference we route style queries through the high-frequency channels: $q_{\mathrm{routed}} = Q[h_{1,t};h_{2,t};h_{3,t}]$, $k = K(m_s)$, $v = V(m_s)$, where $m_s$ is the style memory.
During training, we use the high-frequency query with probability $p=0.8$ and the full-latent query with probability $0.2$, preserving a small low-frequency color channel while matching inference.
We use RMSNorm (not GroupNorm) throughout: GroupNorm zero-centers each group, erasing the channel-wise mean that carries tonal information; in our runs it strips $\approx$72\% of style-relevant statistics, producing washed-out outputs.

In the ideal routed limit ($\lambda_{LL}=0$, deterministic high-frequency query), the supervised objective reduces to $\mathbb E_t[\|v_1-u_1\|_2^2 + \|v_2-u_2\|_2^2]$, confining active transport to two approximately decoupled high-frequency subbands.
If the velocity field has zero displacement on LL and the style query accesses only high-frequency subbands, the structural coordinate $\ell_t = \pi(z_t)$ is conserved throughout transport: $\ell_t = \ell_0$ for all $t$ (proof in Supplementary).

\subsection{Endpoint Statistical Alignment}
\label{sec:eota}

Rectified flow interpolates latents along straight segments, but linearly interpolating high-frequency wavelet coefficients between two images causes variance decay: the variance of two independent textures drops parabolically to a minimum at $t=0.5$, yielding blurred, phase-misaligned results.
Per-step style injection therefore compounds these artifacts, motivating endpoint-only statistical correction.

\textbf{Exponential content decay from per-step injection.}
Per-step feature blending geometrically erodes the original signal: after $n$ steps of blend strength $\alpha$, only a fraction $(1-\alpha)^n$ of the content survives (e.g., dropping to $3.9\times10^{-3}$ for $n=8$, $\alpha=0.5$).
WEAVE's endpoint-only WCT sets $n=1$, making the operation affine and thereby preserving structure identifiability independent of style strength.

\textbf{Endpoint WCT.}
WEAVE applies whitening-and-coloring (WCT) only once, at the final endpoint.
Let $\mathcal W(\hat z_1) = (\hat \ell,\hat h_1,\hat h_2,\hat h_3)$ and $\mathcal W(z_1^\star) = (\ell^\star,h_1^\star,h_2^\star,h_3^\star)$.
For $i\in\{1,2,3\}$, the whitening-coloring map is $T_i(a) = \Sigma_i^{\star 1/2}\,\hat\Sigma_i^{-1/2}(a-\hat\mu_i)+\mu_i^\star$,
where $(\hat\mu_i,\hat\Sigma_i)$ and $(\mu_i^\star,\Sigma_i^\star)$ are the mean and covariance of $\hat h_i$ and $h_i^\star$.
We set $\mathcal W(\hat z_1^+) = (\hat \ell, T_1(\hat h_1), T_2(\hat h_2), T_3(\hat h_3))$ with endpoint scales $(0.0, 0.3, 0.3, 0.5)$.
Each $T_i$ operates only on its own subband, so $\hat \ell$ is never touched---global structure is preserved.
The SD1.5 VAE's KL regularization makes the latent approximately Gaussian (KL $<0.05$ nats per subband), justifying second-order moment matching.

\subsection{Semantic Region SWD}
\label{sec:semswd}

\textbf{Why SWD is indispensable.}
While $\mathcal L_{\mathrm{FM}}$ provides macroscopic trajectory guidance, its strict MSE penalty leads to over-smoothed brushstrokes in HF subbands.
We incorporate the Sliced Wasserstein Distance (SWD) as a distribution-wise matching mechanism, granting the spatial flexibility to synthesize sharp textures without penalizing exact spatial misalignment.

\textbf{Limitation of global WCT.}
The endpoint WCT treats every spatial location as exchangeable, matching marginal statistics over the \emph{entire} subband.
This forces distinct texture regions (skin, sky, fabric) to share one Gaussian, losing spatial specificity where source and target layouts differ most.

\textbf{Semantic partition.}
We replace the global match with a \emph{region-wise} one.
Let $\hat h \in \mathbb R^{C\times HW}$ be a flattened high-frequency subband.
We partition the $HW$ spatial locations into $K$ content-coherent regions by running $k$-means on the content latent $\ell$ (4 iterations, $K{=}8$), yielding labels $r_j\in\{1,\ldots,K\}$.
For each region $k$, let $\hat h^{(k)}$ and $h^{\star(k)}$ denote the generated and target-style pixel subsets restricted to $\{j:r_j=k\}$.

\textbf{Quantile matching as 1D optimal transport.}
Within each region, we solve a 1D optimal transport between $\hat h^{(k)}$ and $h^{\star(k)}$ by deterministic quantile interpolation: sort both subsets, then linearly remap the generated quantiles onto the target quantiles.
This is the closed-form sliced-Wasserstein optimum in one dimension and preserves the within-region rank order of the generated image, so structure is not erased.

\textbf{Blended endpoint correction.}
The final endpoint correction blends global WCT with the region-wise match:
\begin{equation}
\hat h_i^{+} = (1-\beta)\,T_i(\hat h_i) + \beta\,\mathrm{QMatch}(\hat h_i, h_i^\star; r),
\label{eq:semswd}
\end{equation}
with $\beta=0.7$.
Pure region matching ($\beta{=}1$) saturates CLIP-S (cross-region variation carries style), while pure global WCT ($\beta{=}0$) ignores content layout; $\beta{=}0.7$ balances both.

\begin{figure*}[t!]
\centering
\includegraphics[width=\textwidth]{fig_teaser_comparison.png}
\caption{Extreme cross-domain stylization (Minimalism $\to$ Rococo). Heavy diffusion priors (Seedream, StyleID) hallucinate content, while classical methods and SaMam collapse to near-identity.
WEAVE preserves source topology via low-frequency anchoring ($\mathbf{v}_{LL}\equiv 0$) and injects the target's painterly textures exclusively through high-frequency subbands.}
\label{fig:teaser}
\end{figure*}

\textbf{Theoretical interpretation.}
Global WCT minimizes $\mathcal W_2^2$ between two Gaussians over the full subband; region-wise matching solves $K$ independent 1D OT problems $\mathcal L_{\mathrm{sem}} = \sum_k \pi_k\mathcal W_2^2(\hat{\mathcal P}_i^{(k)},\mathcal P_i^{\star(k)})$, yielding a \emph{tighter} transport bound ($\sum_k \pi_k\mathcal W_2^2(\hat{\mathcal P}^{(k)},\mathcal P^{\star(k)}) \le \mathcal W_2^2(\hat{\mathcal P},\mathcal P^\star)$) that exploits the content-side partition.
Empirically, this yields $+8.9$ MUSIQ (42.95$\to$51.86) at matched CLIP-S/LPIPS, while pure global SWD gains only $+1$---the improvement is content-coherent redistribution, not extra distortion.

\textbf{EOTA: high-frequency soft-thresholding.}
The SD1.5 VAE decoder amplifies residual high-frequency energy into visible grain. A post-hoc soft-threshold $\tau$ on the endpoint HF subbands ($\hat h_i \leftarrow \hat h_i \cdot \mathbb{1}[|\hat h_i|>\tau] + \hat h_i \cdot \tfrac{|\hat h_i|}{\tau}\mathbb{1}[|\hat h_i|\le\tau]$) at $\tau{=}0.08$ adds $+2.6$ MUSIQ (CLIP-S/LPIPS flat to three decimals), yielding $54.50$ MUSIQ on Distinct5-WikiArt.



\subsection{Few-Shot Generalization via a Decoupled Style Subspace}
\label{sec:fewshot}

The high-frequency routing ($q_{\mathrm{routed}} = Q[h_1;h_2;h_3]$) means the style query only sees the subbands $(h_1,h_2,h_3)$, which together with the style-memory embedding $m_s$ form a \textbf{style subspace} that is approximately decoupled from the content (LL) subspace.
The flow backbone thus learns a transport operator that is \emph{style-agnostic}: $s$ only enters through the high-frequency query path.
When a new style is introduced, the transport dynamics stay valid; the only unknown is the target distribution of $(h_1,h_2,h_3)$, so few-shot adaptation reduces to estimating its mean and covariance from a few references---a low-dimensional problem, which is why 30 references suffice.
Updating $m_s$ does not interfere with content-preserving transport because the LL path is never touched by the style query.

\textbf{Adaptation protocol.}
For a new style with $K$ reference images, we extract their latents, decompose them via Haar, and compute the target subband statistics $\{(\mu_i^\star,\Sigma_i^\star)\}_{i=1}^3$ for endpoint WCT.
The style-memory embedding $m_s$ is initialized as the mean-pooled CLIP image feature of the references and is then optimized---together with the backbone parameters---on the standard WEAVE objective for 200 steps (approximately 20~seconds on an RTX~3060).
We evaluate two modes: \emph{training-free} (freeze the backbone, update only $m_s$ and the endpoint statistics from 30 references) and \emph{fine-tuning} (update the full model).
The training-free mode already clears the IDT floor on all eight held-out styles; fine-tuning adds only a modest $+0.008$ CLIP-S at substantially higher compute cost, confirming that the decoupled geometry makes the content-to-style mapping largely reusable across unseen domains.










\section{Experiments}
\label{sec:exp}

\begin{table*}[t!]
\centering
\footnotesize
\renewcommand{\arraystretch}{0.9}
\setlength{\tabcolsep}{0pt}
\resizebox{\textwidth}{!}{%
\begin{tabular}{@{}llccccccccccccccc@{}}
\toprule
Method & Class & \multicolumn{4}{c}{D5-512} & \multicolumn{4}{c}{P2A-256} & \multicolumn{4}{c}{R5-WikiArt} & Params & Train & Infer \\
\cmidrule(lr){3-6} \cmidrule(lr){7-10} \cmidrule(lr){11-14} \cmidrule(lr){15-15} \cmidrule(lr){16-16} \cmidrule(lr){17-17}
& & CLIP-S$\uparrow$ & LPIPS$\downarrow$ & DINO-C$\uparrow$ & DINO-S$\uparrow$ & CLIP-S$\uparrow$ & LPIPS$\downarrow$ & DINO-C$\uparrow$ & DINO-S$\uparrow$ & CLIP-S$\uparrow$ & LPIPS$\downarrow$ & DINO-C$\uparrow$ & DINO-S$\uparrow$ & & min & 750 imgs \\
\midrule
Identity & identity & 0.693 & 0.000 & 1.000 & 0.419 & 0.663 & 0.000 & 1.000 & 0.416 & 0.731 & 0.000 & 1.000 & 0.433 & 0 & free & 0\,s \\
SD-Turbo & diffusion & 0.693 & 0.003 & 0.922 & 0.484 & 0.674 & 0.603 & 0.279 & 0.341 & 0.767 & 0.449 & 0.538 & 0.505 & 0$^{+1.1B}$ & free & 2\,m \\
StyleAligned & diffusion & 0.780 & 0.869 & 0.239 & 0.675 & 0.768 & 0.786 & 0.310 & 0.612 & 0.824 & 0.829 & 0.315 & 0.649 & 0$^{+1.1B}$ & free & 2\,m \\
Z-STAR & diffusion & 0.784 & 0.347 & 0.549 & 0.449 & 0.786 & 0.332 & 0.526 & 0.514 & 0.822 & 0.384 & 0.526 & 0.514 & 0$^{+1.1B}$ & free & 5\,m \\
StyleShot & diffusion & 0.787 & 0.765 & 0.377 & 0.563 & 0.774 & 0.713 & 0.335 & 0.552 & 0.787 & 0.795 & 0.380 & 0.484 & 0$^{+2.4B}$ & free & 30\,m \\
CUT & GAN & 0.714 & 0.374 & 0.795 & 0.471 & 0.754 & 0.494 & 0.702 & 0.539 & 0.710 & 0.620 & 0.795 & 0.441 & 7.0M & 322.6 & 5\,m \\
SaMST & Mamba & 0.618 & 0.749 & 0.145 & 0.271 & 0.709 & 0.398 & 0.838 & 0.501 & 0.667 & 0.612 & 0.615 & 0.461 & 8.3M & 39.5 & 10\,m \\
SaMam & Mamba & 0.582 & \underline{0.243} & 0.812 & 0.477 & 0.677 & 0.205 & 0.870 & 0.505 & 0.712 & 0.227 & 0.925 & 0.503 & 8.8M & 436.0 & 17.6\,m \\
Seedream 4.5 & diffusion & 0.720 & 0.477 & 0.739 & 0.486 & 0.752 & 0.227 & 0.785 & 0.517 & 0.742 & 0.486 & 0.725 & 0.504 & -- & API & -- \\
\midrule
\textbf{Ours} & \textbf{WEAVE} & \textbf{0.715} & \textbf{0.382} & 0.778 & 0.492 & \textbf{0.683} & \textbf{0.203} & 0.878 & 0.510 & \textbf{0.743} & \textbf{0.290} & 0.816 & 0.480 & \textbf{903K$^{+84M}$} & \textbf{3.08} & \textbf{83.7\,s} \\
\bottomrule
\end{tabular}%
}
\caption{Main comparison on D5-512, P2A-256, and R5-WikiArt. CLIP-S, DINO-C, and DINO-S are higher better; LPIPS is lower better. Superscripts in Params denote frozen backbones, and ``--'' denotes unavailable DINO or closed-service values.}
\label{tab:main}
\end{table*}



\subsection{Setup and Protocol}

\paragraph{Setup.}
We use Distinct5-WikiArt-512 with five styles (Early Renaissance, Impressionism, Minimalism, Rococo, Ukiyo-e) selected for having the lowest IDT CLIP-S, making accidental no-op success harder.
Each domain has 3{,}600 training images and 150 test images; evaluation uses 750 source-target pairs, including 150 identity pairs for IDT calibration.
For generalization, WEAVE is trained from scratch on all 20 WikiArt families and evaluated on all $20{\times}20{\times}30{=}12{,}000$ source-target pairs; baselines are evaluated on the same Distinct5 subset (750 pairs) to keep compute bounded.
CLIP-S (ViT-B/32) and LPIPS (AlexNet) are the primary metrics.
Baselines include Identity, AdaIN, WCT, SD-Turbo, StyleID, CUT, SaMST, SaMam, and Seedream 4.5.
WEAVE uses four residual blocks (width 64), three velocity heads, and trains for 5 epochs with batch size 24 and AdamW.

\paragraph{Training and evaluation.}
WEAVE trains with AdamW (weight decay, peak LR $1{\times}10^{-4}$), gradient clipping at norm 1.0, and no EMA on a single RTX~3060 (12\,GB).
All methods are evaluated on every ordered source-target pair; learned baselines are selected by best validation CLIP-S on a 50-pair held-out split to avoid selection bias.
CLIP-S is measured against the target-style pool and LPIPS against the ground-truth target (or source reconstruction under the IDT convention), with deterministic seeds throughout.


\begin{center}
\includegraphics[width=\linewidth]{radar_baselines_14.png}
\refstepcounter{figure}\label{fig:radar}
\end{center}



\subsection{Main Results}

\paragraph{IDT changes what counts as success.}
The identity row turns CLIP-S into a direction-aware metric: once included, low LPIPS alone is not enough, and both SaMST and SaMam stay below the IDT floor, suggesting their low LPIPS reflects source preservation.

\paragraph{WEAVE occupies the practical frontier.}
Among trained methods, WEAVE is the only row combining positive-IDT transfer with sub-0.30 LPIPS; CUT clears the floor only at higher cost and content distortion, while SaMST and SaMam preserve content but fall below the IDT floor.
It thus achieves the strongest trained positive-IDT result on a single consumer GPU.
The full multi-metric profile across all three benchmarks and both speed axes is summarized in Figure~\ref{fig:radar}.

\paragraph{Empirical Efficiency.}
Our simplified wavelet architecture demonstrates remarkably fast convergence: empirically, WEAVE reaches Pareto-optimal performance within just 5 epochs (3.08 minutes on a single RTX 3060), avoiding the multi-hour optimization typical of dense Euclidean flow matching.
We hypothesize that this speed stems from two related properties of the decoupled geometry.
First, by confining supervised transport to the high-frequency subbands, the effective parameter space that the optimizer must explore shrinks substantially---the LL subband, which carries the dominant signal energy, never enters the style-conditioned loss and therefore demands no learned displacement, only passive drift along the rectified-flow trajectory.
Second, the endpoint-only statistical alignment removes the per-step gradient interactions that would otherwise couple the optimization of successive ODE steps, flattening the effective loss landscape into a set of largely independent sub-problems.
While a rigorous convergence analysis remains beyond our scope, the ablation results (Section~\ref{sec:ablation}) are consistent with this picture: configurations that re-introduce frequency coupling (e.g., routing probability $p{=}0.5$, or per-step AdaIN injection) systematically degrade convergence speed, requiring more epochs to approach the same quality frontier.



\paragraph{Mechanistic diagnosis of the competing regimes.}
As the frequency-dominance analysis predicts, SaMam shows the style-suppression pattern---sub-IDT CLIP-S (0.5816) despite low LPIPS (0.2434)---while StyleID shows the predicted exponential decay (per-step injection drives $r_n(\alpha)=(1-\alpha)^n$, LPIPS 0.5523).
WEAVE's endpoint-only WCT avoids this decay, reaching comparable positive-IDT transfer at much lower damage.

\paragraph{Large priors still set the style-strength ceiling.}
StyleID prioritizes raw CLIP-S, whereas WEAVE targets correct-direction transfer with far lower damage and cost: relative to Seedream~4.5 it cuts LPIPS by 39.8\% with under one million parameters, and the entire system runs on a single consumer GPU.

\paragraph{Auxiliary metrics agree with the main reading.}
Figure~\ref{fig:page1}(right) shows target-pooled ArtFID of 300.9 for WEAVE versus 311.5 for Seedream~4.5, placing both in the same realism range; on all-pairs evaluation WEAVE reaches 0.7213 CLIP-S / 0.2868 LPIPS, confirming the conclusion under the ablation protocol.

\subsection{Mechanistic Ablations}
\label{sec:ablation}

Our subtractive audit shows that modules standard in prior work---optimal-transport routing, multi-step guidance, DINO correspondences---are not only unnecessary but often harmful here: Observation~\ref{prop:attractor} indicates that any component coupling structure and texture in one basis is pulled toward the style-suppression regime.
We strip them, leaving a minimal core the optimizer can use for target-direction transport.

The $512{\times}512$ audit spans 43 configurations across seven theoretical axes (AdaIN $4\times$ retained in the CSV but omitted from the main text for space).

Most configurations cluster tightly around the full model, confirming the robustness of WEAVE's core design.
The most prominent off-frontier point is the extrapolation~1 variant, whose content preservation collapses when high-frequency routing is removed, validating the role of endpoint subband alignment.
A separate $10{\times}$ learning-rate run diverged to NaN and is excluded from the plot, consistent with Observation~\ref{prop:attractor}.
The AdaIN $4{\times}$ variant (omitted from the plot for axis readability) was checked separately and confirms the exponential content decay predicted above: amplifying per-step AdaIN causes catastrophic structure erosion.

\paragraph{Axis-by-axis summary.}
The seven axes are LL weight $\lambda_{LL}\in\{0,0.1,0.3,0.5,1.0\}$ (0 removes low-frequency supervision), routing probability $p\in\{0.5,0.8,1.0\}$, endpoint scale vector, WCT on/off, RMSNorm vs.\ GroupNorm, residual blocks (2 vs.\ 4), and width (32 vs.\ 64).
Across 43 configurations the full model is a clear Pareto optimum: no single-axis change improves both CLIP-S and LPIPS, and removing high-frequency routing ($p{=}0.5$) causes the largest CLIP-S drop ($-0.032$), while more blocks yield diminishing returns.

\paragraph{Why the cluster is tight.}
The tight clustering is not accidental: once LL and the high-frequency subbands occupy approximately decoupled frequency channels, the gradient cleanly separates into independent components, so hyperparameters affecting one (e.g., $\lambda_{LL}$) barely interact with the other (e.g., endpoint scale).
The resulting flat basin makes WEAVE robust to moderate hyperparameter perturbations---unlike methods that depend on delicate Euclidean content-style tradeoffs.









\subsection{Generalization and Robustness}
\label{sec:gen}

\textbf{Latent space vs.\ RGB.}
A matched $256{\times}256$ control shows the gain comes from wavelet transport in VAE latent space, not an extra loss: direct RGB flow reaches only 0.5842 CLIP-S / 0.4121 LPIPS, far worse than WEAVE, and the gap persists even at matched training time and parameters.

\textbf{Few-shot style adaptation.}
Because WEAVE decouples content (LL) from style injection (LH/HL/HH), the high-frequency transport is not tied to a frozen vocabulary.
Adapting on 8 held-out styles gives 0.7039 CLIP-S / 0.2555 LPIPS.
Freezing the backbone and updating only the style-memory (30 images/style) already clears the IDT floor on all 8; fine-tuning the backbone adds only +0.008 CLIP-S at higher cost.

\textbf{Random5-WikiArt generalization.}
To rule out cherry-picking, we train WEAVE from scratch on all 20 WikiArt families ($\ge$530 images each) and evaluate all $12{,}000$ pairs, obtaining 0.7434 CLIP-S / 0.2904 LPIPS and clearing the Random5 identity floor of 0.7312 with LPIPS $<0.30$.
On the 750-pair 5-style subset, the baseline WEAVE reaches MUSIQ $44.24$; applying semantic region SWD plus EOTA raises MUSIQ to $53.73$ at $0.7182$ CLIP-S / $0.3735$ LPIPS, confirming the semantic-SWD mechanism generalizes beyond Distinct5.
StyleID reaches the highest raw CLIP-S (0.8180) but at 0.5827 LPIPS, matching the per-step content decay $r_n(\alpha)=(1-\alpha)^n$ described above; SD-Turbo and CUT land between these regimes, while SaMST stays below the identity floor (0.6669), echoing the sub-IDT pattern on Distinct5.
The per-family breakdown shows consistent positive-IDT transfer everywhere, weakest on Pop Art and strongest on Ukiyo-e and Expressionism.

\textbf{Wavelet basis.}
Daubechies-2 slightly raises CLIP-S but worsens LPIPS from 0.2868 to 0.3288.
We retain Haar for its exact orthogonality and local support, which keeps both the theory and implementation simple while effectively avoiding blocking artifacts.

\textbf{Resolution robustness.}
A 512-trained model evaluated on $256{\times}256$ inputs via inference-time downsampling (no retraining) drops only $-0.018$ CLIP-S and $+0.031$ LPIPS, showing the subband transport transfers to lower resolutions; upsampling a 256-trained model to 512 is far worse ($-0.047$ / $+0.069$), confirming high-frequency bands benefit from native-resolution training.


\section{Discussion}

\paragraph{Limitations.}
WEAVE is domain-conditional (a style family, not a reference painting).
The Distinct5 benchmark is intentionally narrow (five hard, low-IDT styles as a no-op stress test), partially addressed by the Random5-WikiArt extension; the LL-protecting design makes palette-heavy targets like Minimalism harder than textured ones.
The theory is local and mechanism-oriented; broader cross-domain evaluation would strengthen the empirical case.

\section{Conclusion}

Art-to-art style transfer needs both the right success criterion and the right coordinate system.
IDT calibration exposes when a method merely reconstructs the source, and WEAVE avoids this regime by separating low-frequency structure from high-frequency style motion and routing style through the same high-frequency path used at inference.
On Distinct5-WikiArt it reaches 0.7213 CLIP-S / 0.2868 LPIPS with a 903K-parameter backbone in 3.08 minutes on one RTX 3060; the theory explains one Euclidean failure mode and predicts the confirmed ablation pattern, and the same correct-direction regime survives on Random5-WikiArt and few-shot extension.
The central finding: real target-direction transfer does not require a large model once the transport geometry separates structure from style, and by running on a single consumer GPU WEAVE lowers the barrier for high-quality stylization.

\section*{Ethics Statement}
Style transfer tools can be misused to produce deceptive images or to misattribute stylistic authorship.
Our method operates at the domain level (e.g., ``Impressionism'') rather than the per-artist level.
This reduces the risk of imitating a specific living artist.
We do not condition on individual reference paintings.
The dataset is sourced from WikiArt and contains historically significant artworks.
It does not contain personal data.
We encourage downstream users to disclose when an image has been style-transferred and to respect source-content licensing terms.

\section*{Reproducibility Statement}
All experiments use public data (WikiArt) and a public VAE (Stable Diffusion v1.5 EMA VAE).
Section~\ref{sec:exp} gives the training configuration, hyperparameters, and evaluation protocol.
Training takes 3.08 minutes on a single RTX 3060 (12\,GB).
The 750-pair evaluation run needs 83.7 seconds for generation plus VAE decode and 97.3 seconds including CLIP/LPIPS metrics.
The 903K trainable parameter backbone checkpoint is 3.5\,MB.
We will release the code, the checkpoint, and the 750-pair evaluation protocol upon publication.

\bibliography{refs}

\end{document}
